% !TeX program = xelatex
% Compile twice with XeLaTeX. All body text, tables and bibliography are embedded.
% Based on geshi/template/ydzjreport.cls, unmodified v1.3.
\documentclass[stage2,UTF8,fontset=fandol]{ydzjreport}
\usepackage{bm,longtable,subcaption,etoolbox}
\setlength{\headheight}{15pt}
\setcounter{secnumdepth}{3}
\setcounter{tocdepth}{2}
\AtBeginEnvironment{tabularx}{\linespread{1.1}\selectfont}
\AtBeginEnvironment{longtable}{\linespread{1.1}\selectfont}
\AtBeginEnvironment{thebibliography}{\raggedright}
\pretocmd{\section}{\Needspace{14\baselineskip}}{}{}
\pretocmd{\subsection}{\Needspace{10\baselineskip}}{}{}
\pretocmd{\subsubsection}{\Needspace{5\baselineskip}}{}{}
\titleformat{\paragraph}[runin]{\normalsize\heiti\color{ydzjblueiii}}{}{0pt}{}
\titlespacing{\paragraph}{0pt}{5pt}{1em}
\DeclareRobustCommand{\TBD}[1]{{\color{ydzjgrey}\textbf{【待补充：#1】}}}
\newcommand{\pending}[1]{\par\noindent\textbf{证据补充与适用边界：}#1\par}
\newcommand{\planned}[1]{\par\noindent\textbf{后续拟开展：}#1\par}
\newcommand{\placeholderbox}[2]{%
  \begingroup\setlength{\fboxsep}{7pt}%
  \fbox{\begin{minipage}[c][#2][c]{\dimexpr\linewidth-2\fboxsep-2\fboxrule\relax}%
  \centering\small\textbf{真实图片预留位置}\par\medskip #1%
  \end{minipage}}\endgroup}
\newcommand{\placeholderfigure}[4]{%
  \begin{figure}[!htbp]\centering
  \placeholderbox{#3}{#4}\caption{#2}\label{#1}\end{figure}}
\newcommand{\placeholderpair}[7]{%
  \begin{figure}[!htbp]\centering
  \begin{subfigure}[t]{0.48\linewidth}\centering
    \placeholderbox{#4}{#7}\caption{#3}\label{#1:a}
  \end{subfigure}\hfill
  \begin{subfigure}[t]{0.48\linewidth}\centering
    \placeholderbox{#6}{#7}\caption{#5}\label{#1:b}
  \end{subfigure}
  \caption{#2}\label{#1}\end{figure}}
\newcommand{\reportfigure}[4]{\begin{figure}[!htbp]\centering
  \includegraphics[width=#4]{#3}\caption{#2}\label{#1}\end{figure}}
\stagetwosetup
  {中期汇报}
  {项目团队}
  {微通道结构设计与动态调控}
  {报告起止日期待确认}
  {非均匀热负荷；区域化连接；协同调控；物理融合代理模型}
  {已形成局部结构验证与最小闭环基础，正在完善动态模型、数据集和系统级证据。}
\setshorttitle{微通道热管理 · 中期研究进展}
\hypersetup{pdftitle={基于时空非均匀性的微通道热管理系统的结构设计与动态调控策略——阶段二研究进展}}
\begin{document}
\makestagetwohead
\begin{keybox}
\textbf{项目名称：}基于时空非均匀性的微通道热管理系统的结构设计与动态调控策略。
\end{keybox}

\par\medskip
\noindent\textbf{项目概述}\par
时空非均匀热负荷使微通道冷板同时面临局部热点抑制、温度分布协调和泵送能耗控制等要求。针对均匀结构与固定流量运行方式可能存在的供需不匹配，本项目初期提出区域化结构设计及单元连接整流、非均匀动态热负荷下流量分配与协同控制两条并行路径。前者着重改善不同热负荷区域之间的结构适配及连接处的流动组织，后者着重通过双泵和行级出口阀门协调冷却资源。随着局部结构参数和工况组合不断扩展，高保真计算成本成为新的方法问题，项目进一步形成物理机理融合的流热性能代理模型与计算加速方向，逐步构建结构、控制与计算方法相互支撑的研究体系。

截至中期，团队已开展代表性翼型扰流结构建模与数值仿真、参数化和自动化流程探索、多轮网格敏感性检验，以及温度相关水物性 UDF 的修正与对照；已经形成面向局部结构研究的数值工作基础。在区域化连接结构方面，已确定连接方案，依据团队阶段进展说明，已有数值仿真对照验证其性能优势，相关研究接近阶段性完成，当前重点为对照数据归档、最终性能评价与机理解释。动态调控方面，已形成十行单向流路、双泵和行级出口阀门的基本架构，开展随机变化热源模拟并完成最小闭环的初步可行性验证。代理模型方向已确定技术路线，正在推进高保真样本建设、物理尺度分析及基础算法研究，尚未形成经独立测试验证的模型精度、泛化能力或计算加速结果。

本报告依照项目真实演化过程梳理问题、方法与阶段证据，并区分局部连接验证、装配几何准备及整板系统验证的不同层级。现有翼型网格资料表明不同指标对加密的敏感性并不一致，物性实现对照反映的是实现一致性而非全部流热模型的绝对精度。下一阶段将完成连接结构成果整理，深化考虑能量守恒、热阻热容及泵阀特性的动态模型，并以统一数据口径评价物理融合代理模型的数据效率和累计计算成本，最终开展面向温度响应、温度均匀性和泵送代价的系统级协同验证。

\section{本月研究任务}
\subsection{任务目标}
本报告覆盖截至中期的累计研究进展；具体月度起止日期尚待确认。总体任务是完成区域化连接结构的阶段性验证、建立动态反馈研究基础，并为高保真数据驱动的快速预测方法建立可靠数值平台。
\subsection{完成情况}
见表~\ref{tab:monthly_status}。
\begin{table}[htbp]\centering\small
\caption{阶段研究任务完成情况}\label{tab:monthly_status}
\begin{tabularx}{\linewidth}{@{}Y c Y@{}}\toprule\rowcolor{ydzjlight} 
\rowcolor{ydzjlight} 任务 & 状态 & 说明 \\\midrule
区域化连接结构设计与验证 & 部分完成 & 已确定方案并有局部对照；最终成对结果和机理归档待完成。\\
动态热源与最小反馈链条 & 已完成 & 仅指初步可行性；完整日志与系统性能评价仍待补。\\
数值平台与物性实现核对 & 部分完成 & 已有多轮检验；生产网格与新工况准入尚待确认。\\
物理融合代理模型研究 & 部分完成 & 正在建库和研究基础算法，尚无独立测试精度。\\
整板流热与协同系统验证 & 未开始 & 当前整板仅完成几何和热源映射，不能替代流热求解。\\
\bottomrule\end{tabularx}\end{table}
\subsection{研究背景与项目目标}
\label{ch:background}
\subsubsection{时空非均匀热负荷下的散热需求}
微通道利用较小的流动尺度和较大的固液接触面积实现紧凑冷却，是高热流密度热管理的重要研究对象。早期微通道热沉研究为这一技术路径提供了基础\cite{tuckerman1981}。对于本项目关注的冷板，热源并非在整个受热面上均匀分布：局部区域可能长期维持高负荷，另一些区域处于低负荷，而热点位置和强度还可能随时间变化。因此，平均热流密度不足以代表系统最困难的散热条件，需要同时描述空间分布、波动幅度、变化时间尺度及持续时间。

空间非均匀性意味着冷却需求存在区域差异。若全部区域采用同一种微通道结构和相同流量，低负荷区域可能获得超过需求的冷却资源，高负荷区域则仍受局部热阻限制。时间非均匀性进一步要求冷却供给随负荷变化调整，但固体热容、流体输运时延和泵阀响应决定了温度不会立即随控制指令变化。项目因此将热管理问题视为结构适配与动态资源分配的共同问题。图\ref{fig:01}用于说明这种空间分布与时间响应的关系。
\reportfigure{fig:01}{时空非均匀热负荷下的微通道热管理问题}{figures/diagrams/fig_01.pdf}{0.98\linewidth}

\subsubsection{微通道热管理面临的问题}
局部强化结构能够改变流动路径、增加有效换热面积或促进热边界层发展，但也可能增加摩擦与局部阻力。评价结构时，应同时考虑最高温度、温度差异、压降和泵送功率，而不能仅依据某一温度指标判断优劣。对于模块化冷板，单元内部性能还会受到连接段的影响：连接处流动收缩、扩张、转向与速度分布重构可能改变下游单元的入口条件，使单元性能难以直接线性叠加。

固定流量方法实施简单，却缺乏对热点迁移和热负荷波动的主动响应。增加总流量可能降低部分区域温度，但其能耗代价和区域分配效果需要联合评价。与此同时，复杂几何、温度相关物性以及瞬态边界增加了高保真计算成本。上述问题并非彼此独立：结构改变水力阻力，水力阻力影响控制分配，控制运行又改变结构的实际流热工作点。

\subsubsection{项目的原始研究思路}
项目最初围绕两条路径展开。第一条路径将冷板划分为具有明确边界的基本单元，使局部结构与代表性热负荷相适应，并通过连接整流改善单元间流动组织。第二条路径面向热负荷随时间变化的情况，研究利用区域温度反馈协调泵和阀门的流量分配。两条路径自初期即构成整体框架，前者提供被控对象，后者调节其运行状态，并非先后互相替代的方案。

这种设计逻辑也限定了研究的证据要求：结构方向需提供一致工况下的基准对照，控制方向需提供热源扰动与温度响应之间的闭环记录。代理模型并不是最初与两条路径完全平行设定的独立目标，而是在研究规模扩大后，为支撑重复性能预测而形成的方法拓展。

\subsubsection{项目总体研究目标}
项目的总体目标是在时空非均匀负荷条件下，形成从局部结构、单元连接到动态运行的可验证研究链条。中期考核的重点是说明这一链条已推进到何种程度，而非提前宣布系统最优或所有任务完成。表\ref{tab:01}给出目标与三个方向之间的对应关系。
\begin{table}[htbp]
\centering\small
\caption{研究目标、对应方向与中期状态}\label{tab:01}
\begin{tabularx}{\linewidth}{lYY}
\toprule
\rowcolor{ydzjlight} 方向 & 研究目标 & 阶段依据与交付边界 \\
\midrule
方向二 & 区域结构适配与连接性能验证 & 方案确定；已有局部对照支持优势；最终量化与机理待归档。 \\
方向三 & 动态温度响应与泵阀协同 & 最小闭环初步可行；物理模型与多算法比较正在推进。 \\
方向一 & 物理融合预测与计算加速 & 已有数值与自动化基础；数据集与算法研究正在推进。 \\
系统协同 & 统一对象下评价温度、均匀性和代价 & 整板几何与热源映射已有记录；系统流热与控制验证拟开展。 \\
\bottomrule
\end{tabularx}
\end{table}

区域化连接方向以可复现的性能对照与机制解释为阶段交付；动态调控方向以可行的反馈链条和逐步深化的物理模型为阶段交付；代理模型方向以可靠数据、算法比较及成本评价为后续交付。团队共四名本科生，实际姓名、职责与周期在报告头和任务表中保留待补项，建议分工不作为已经发生的贡献记录。
\subsection{项目总体技术路线与研究演化}
\label{ch:evolution}
\subsubsection{项目初期的研究框架}
项目的出发点是冷却需求与冷却供给在空间和时间上的不匹配。初期将这一问题分解为结构层面的区域适配和运行层面的动态分配：结构路径讨论不同区域需要怎样的通道以及怎样连接，控制路径讨论既有冷板需要在什么时刻向哪些区域提供多少流量。图\ref{fig:02}用于保留初始双路径框架，避免在报告中将后续衍生方向误写成初始设定。
\reportfigure{fig:02}{项目初始的结构设计与动态调控双路径}{figures/diagrams/fig_02.pdf}{0.98\linewidth}
初期框架的价值在于建立统一研究对象。结构研究并非只服务于稳态散热，控制研究也不是脱离通道阻力的独立算法演示。连接段和局部结构产生的压降特性，将成为控制系统分配流量时必须考虑的条件；动态热源则为判断区域化设计是否具有实际适用性提供新的验证场景。

\subsubsection{区域化结构与动态控制的并行推进}
在结构路径中，团队完成基本单元划分、代表性负荷分类和连接结构方案设计，并推进局部数值对照。已有阶段说明表明连接方案已确定且性能优势获得仿真支持，研究重点由方案生成逐渐转向结果收敛、对照归档与机理解释。在控制路径中，团队形成双泵与行级出口阀门方案，开展随机热负荷模拟，并完成最小闭环的初步验证，重点由基本反馈可行性逐渐转向物理模型和控制性能评价。

两条路径在推进深度上有所不同，这种差异是中期状态的实际体现。连接结构接近阶段性完成，不等于整板与控制系统均已验证；最小闭环实现也不等于多输入耦合系统已经完成全面控制比较。报告据此分别设置结构成果与控制深化部分，避免用统一的“已完成研究”概括不同成熟度。

\subsubsection{高保真CFD计算需求与代理模型方向的提出}
随着翼型扰流单元、连接区域和不同热边界的研究展开，几何参数与工况的组合不断增加。同一个候选结构往往需要经历网格检验、物性核对以及多个边界工况计算，单次求解以外还存在几何处理、网格生成、监测与数据提取等成本。这使重复调用高保真 CFD 的代价逐渐成为限制研究扩展的因素。

团队由此提出新的方法问题：能否在明确适用范围和误差代价的前提下，用物理建模与机器学习减少重复求解？这一问题促成方向一的形成。其核心是建立高效预测方法并研究物理信息如何提高模型可靠性，几何参数化主要服务于样本生成与泛化检验，不以寻找最优几何作为本方向的中心成果。图\ref{fig:03}展示工程研究如何产生计算方法需求，而不采用三个方向从一开始并列的叙事。
\reportfigure{fig:03}{由工程问题驱动的项目真实研究演化}{figures/diagrams/fig_03.pdf}{0.90\linewidth}

\subsubsection{三个研究方向的最终技术体系}
当前体系由三个方向组成：方向一为物理机理融合的微通道流热性能代理模型与计算加速；方向二为区域化冷板连接结构优化与性能验证；方向三为考虑流动传热物理机理的非均匀动态负荷协同调控。方向二提供局部结构、连接设计和性能对照，方向三提供系统运行需求和动态场景，方向一利用可靠 CFD 数据形成高效预测工具，并将误差与适用边界反馈给前两者。

这种协同关系仍需分阶段实现。代理模型在通过独立测试前不能直接替代高保真评价；控制模型在完成泵阀标定前不能保证准确流量分配；局部连接改善也不能自动推广为整板优势。图\ref{fig:04}概括对应数据、参数和验证需求的流向，区分当前已有支撑与后续拟建立的联动。
\enlargethispage{\baselineskip}
\reportfigure{fig:04}{三个研究方向的数据与模型协同关系}{figures/diagrams/fig_04.pdf}{0.90\linewidth}
\section{进阶阅读与理论支撑}
\subsection{本阶段文献清单}
本节沿用原报告已核实的三项文献，不从格式范例引入其他项目的仿真数据或参考结论。文献与用途对应关系见表~\ref{tab:reading}。
\begin{table}[htbp]\centering\small
\caption{已有文献与本项目的方法对应}\label{tab:reading}
\begin{tabularx}{\linewidth}{@{}Y Y Y@{}}\toprule\rowcolor{ydzjlight} 
\rowcolor{ydzjlight} 文献 & 对应问题 & 本项目采用范围 \\\midrule
Tuckerman 与 Pease（1981）\cite{tuckerman1981} & 紧凑微通道散热的技术背景 & 作为背景，不代替本项目结构对照。\\
Chen 与 Guestrin（2016）\cite{chen2016} & 表格型性能数据的基础回归 & 拟作为基础算法，不预设预测精度。\\
Raissi 等（2017）\cite{raissi2017} & 将物理规律引入学习过程 & 作为方法参照，结合输出形式选择特征、残差或约束路线。\\
\bottomrule\end{tabularx}\end{table}
\subsection{对本方向的启示}
现有方法依据支持先建立可信数值样本和基础预测，再检验物理信息的增益。流热守恒关系可用于特征设计、一致性检查或简化基线，但不能直接确定局部最高壁温；树模型也不能照搬神经网络的偏微分方程自动微分约束。因此，本项目拟先比较基础模型、物理特征增强及残差学习，再评价适用条件明确的约束网络，并将独立测试、消融与累计成本作为判断依据。相关推导和适用边界保留在研究进展中。
\section{研究进展}
\subsection{已完成的工作}

表\ref{tab:15}以成熟度和证据状态总结当前进展。方向二已有确定方案与局部仿真验证，是目前最接近阶段收尾的方向；方向三已实现基本反馈可行性，正在深化物理模型；方向一已建立部分数值与流程基础，处于数据集与算法研究阶段。图\ref{fig:29}概括已有节点、正在推进任务与成果边界，不使用无依据的完成百分比。
\begin{table}[htbp]
\centering\small
\caption{三个方向当前进展与证据边界}\label{tab:15}
\begin{tabularx}{\linewidth}{lYY}
\toprule
\rowcolor{ydzjlight} 方向 & 已完成或已开展 & 正在开展与后续 \\
\midrule
方向二 & 分区、连接方案确定、局部数值对照支持优势 & 最终量化归档、关键参数与机制；接近阶段完成。 \\
方向三 & 动态框架、双泵行阀方案、热源模拟、最小闭环初步验证 & 灰箱模型、泵阀特性、控制比较；完整系统效果未形成。 \\
方向一 & 翼型数值研究、自动化探索、网格检验、物性修正、技术路线 & 样本集、尺度推导、基础及融合算法；精度与加速未产生。 \\
共享平台 & 当前整板几何、热源映射；局部方法与指标核验 & 整板流体域、网格、求解及系统协同验证。 \\
\bottomrule
\end{tabularx}
\par\smallskip\begin{minipage}{\linewidth}\footnotesize 连接优势和最小闭环阶段事实来自团队说明，原始成对数据与闭环日志待补；不以文件数量计算完成率。\end{minipage}
\end{table}

\reportfigure{fig:29}{三个研究方向的阶段性进展示意}{figures/diagrams/fig_29.pdf}{0.98\linewidth}


第一，形成了以空间适配与时间调控为起点的研究框架，并在仿真需求推动下拓展出计算加速方法方向。其价值在于明确了工程对象、动态需求与计算方法之间的关系，使后续研究能够共同服务于同一热管理问题。第二，完成区域化基本单元、连接方案及局部验证，取得了团队已确认的性能优势，下一步是将其转化为数据与机理完整的阶段结果。

第三，开展翼型模型、网格敏感性和压力统计核验，修正温度相关水物性并进行实现对照，使数值工作从几何生成逐步进入评价口径与可靠性控制。第四，建立动态热源与反馈控制的初步研究基础，验证最小闭环的可行性，为后续物理灰箱与算法比较提供起点。第五，已有当前整板装配与热源映射验证，支持结构与空间热源的统一配置，但整板流热性能仍须后续求解。

上述成果分别对应研究框架、局部结构、数值方法、控制原型和系统几何准备，证据层级不同。已有资料中的局部物性实现计时属于数值实现改进，网格检验属于可靠性评价，几何装配属于工程准备，均不能替代代理模型精度或系统控制收益。当前没有已提供的论文发表、专利或实测节能成果，本报告不填写此类状态。

\subsection{关键结果}
已有量化记录见表~\ref{tab:key_results}；具体统计口径、原始数值和适用边界见后续数值平台与结构部分。以下数据均来自原报告，不表示新增求解结果。
\begin{table}[htbp]\centering\small
\caption{已有量化记录与证据边界}\label{tab:key_results}
\begin{tabularx}{\linewidth}{@{}Y Y Y@{}}\toprule\rowcolor{ydzjlight} 
\rowcolor{ydzjlight} 记录 & 量化信息 & 证据边界 \\\midrule
网格敏感性 & 7 档，260169--4610769 单元，统计时刻 0.2 s & 最细档时间步历史有差异，全序列未通过联合判据。\\
温度相关水物性 & 280--360 K，101325 Pa & 拟合适用范围，不表示任意压力与温度均有效。\\
物性实现对照 & 共同 0.137 s 加热面最高温差约 $3.197\times10^{-6}$ K & 实现一致性核对，不代表模型绝对精度。\\
整板几何 & $40\,\mathrm{mm}\times40\,\mathrm{mm}$，$10\times10$ 单元 & 几何与热源映射已完成，整板流热待验证。\\
静态热源布局 & 低/中/高负荷为 70/27/3 格 & 整格加热 2640 W；仅主体加热 2021.25 W。\\
\bottomrule\end{tabularx}\end{table}
\subsection{与预期的差异}
\textbf{网格检验方面：}加密本应减小离散误差，但现有七档序列并未在平均温度、最高温度与压差上共同单调收敛，最细档还混有不同时间步历史。因此，不能直接宣布网格无关，下一步需统一初场和时间推进后复核。

\textbf{研究路线方面：}原始结构与控制双路径在推进中产生了重复高保真计算成本问题，因此进一步形成代理模型方向。新增方向目前只有数值与流程基础，尚未达到可量化预测精度或累计加速收益的阶段。

\textbf{验证层级方面：}局部连接方案已有团队确认的定性优势，当前整板则仅完成几何准备；两者成熟度不同。后续必须分别完成局部成对结果归档和整板流热求解，不能将局部优势直接推广为系统收益。
\subsection{数值研究平台与研究方法}
\label{ch:methods}
\subsubsection{几何建模与计算域}
数值研究以整体冷板、局部单元和连接区域三个层级组织。整体设计为 $40\,\mathrm{mm}\times40\,\mathrm{mm}$ 冷板，按 $10\times10$ 个基本单元划分，每个单元节距为 $4\,\mathrm{mm}\times4\,\mathrm{mm}$。节距包含结构主体及连接或间隔空间，不能将其与主体流道的有效尺寸混同。已读取的当前整板几何说明给出 $3.5\,\mathrm{mm}\times3.5\,\mathrm{mm}$ 主体、$3\,\mathrm{mm}$ 总厚度及十条并排流路；这些数据对应具体版本，并非所有历史模型的统一尺寸。表\ref{tab:02}汇总已知参数与待确认尺寸。
\begin{table}[htbp]
\centering\small
\caption{冷板及代表性微通道几何参数}\label{tab:02}
\begin{tabularx}{\linewidth}{lYY}
\toprule
\rowcolor{ydzjlight} 参数 & 已知值或待补值 & 适用版本与说明 \\
\midrule
整体外形 & $40\times40\times3\,\mathrm{mm}$ & 当前whole\_plate装配说明；厚度不外推至全部历史模型。 \\
单元划分 & $10\times10$，100格 & 用户进展说明及当前热源布局。 \\
单元节距 & $4\times4\,\mathrm{mm}$ & 包含主体和接口/间隔。 \\
主体平面尺寸 & $3.5\times3.5\,\mathrm{mm}$ & 当前整板装配。 \\
行内连接长 & $0.5\,\mathrm{mm}$ & 当前整板及40 mm条带装配。 \\
端部连接长 & $0.25\,\mathrm{mm}$ & 入口与出口端分别设置。 \\
代表翼型 & NACA2416，0.45 mm弦长 & 当前40 mm条带代表配置；非最优结构。 \\
翼型排布 & 9/8/9/8/9，43柱，0度迎角 & 该代表配置；其他样本单独登记。 \\
流道与最小间隙 & \TBD{主体截面、柱高、最小间隙} & 逐一绑定几何文件及版本。 \\
\bottomrule
\end{tabularx}
\end{table}

图\ref{fig:05}和图\ref{fig:06}分别预留整体布局和局部结构图。翼型模型以 NACA2416 等结构为代表；局部截面、扰流柱和连接段应在图中注明坐标、尺寸及版本。固体域、流体域及交界面需要作为共轭传热对象共同处理，流体连通性和固液接触关系均须检查。当前整板已有几何与热源排布验证，仍需后续流体域处理、网格与求解，不能引用局部模型的结果作为整板计算结果。
\placeholderfigure{fig:05}{整体冷板分区与基本单元几何}{40×40 mm外形、10×10网格、4 mm节距、3.5 mm主体、连接与条间填充；可采用现有whole\_plate截面图补充尺寸。}{33mm}
\placeholderpair{fig:06}{代表性局部微通道结构模型}{局部主体}{补入翼型扰流或其他已核实结构，注明流向、尺度、版本。}{连接区域}{补入确定方案及连接长度，说明与主体接口。}{30mm}

\subsubsection{CFD数值模型与基本控制方程}
以下方程是描述单相流动与共轭传热的通用物理基础，用于说明变量和守恒关系；它们不意味着项目已采用全部高阶物理项或所有模型假设均已验证。已读取的代表性物性对照案例采用三维、层流、能量方程及温度相关水物性；整板规划沿用层流和共轭传热思路。不同工况的稳态或瞬态选择、离散格式和求解容差应逐案例确认，表\ref{tab:03}保留具体配置。

对密度可能随温度变化的单相流体，质量守恒写为
\begin{equation}\frac{\partial\rho}{\partial t}+\nabla\cdot(\rho\bm u)=0.\label{eq:mass}\end{equation}
其中 $\rho$ 为密度，$\bm u$ 为速度，$t$ 为时间。只有在恒密度近似下才可简化为 $\nabla\cdot\bm u=0$。温度相关密度启用后，瞬态质量诊断应计入储存项。
\begin{equation}\begin{aligned}
\frac{\partial(\rho\bm u)}{\partial t}+\nabla\cdot(\rho\bm u\otimes\bm u)
&=-\nabla p+\nabla\cdot\bm\tau+\rho\bm g,\\
\bm\tau&=\mu\left[\nabla\bm u+(\nabla\bm u)^{\mathsf T}-\frac23(\nabla\cdot\bm u)\bm I\right].
\end{aligned}\label{eq:momentum}
\end{equation}
$p$ 为压力，$\mu$ 为动力黏度，$\bm g$ 为重力加速度，$\bm I$ 为单位张量。该式适用于牛顿流体、采用 Stokes 假设的黏性应力表述；代表性案例关闭重力。层流模型能否覆盖未来参数范围，应结合雷诺数与流动现象重新判断。

在低马赫数、忽略黏性耗散与压力功、无相变条件下，流体显热方程可表为
\begin{equation}
\frac{\partial(\rho h_f)}{\partial t}+\nabla\cdot(\rho\bm u h_f)
=\nabla\cdot(k_f\nabla T_f)+s_f,
\quad h_f(T)=\int_{T_0}^{T}c_p(\theta)\,\mathrm d\theta.\label{eq:fluidenergy}
\end{equation}
$T_f$ 为流体温度，$k_f$ 为导热系数，$c_p$ 为定压比热，$s_f$ 为体积热源，$T_0$ 为焓参考温度。表面热流边界与体积热源的单位不同，不能混用。温度相关比热时，显热应由积分获得，不宜始终采用常数 $c_p\Delta T$。
\begin{equation}
\rho_s c_s\frac{\partial T_s}{\partial t}=\nabla\cdot(k_s\nabla T_s)+s_s.\label{eq:solidenergy}
\end{equation}
该式描述固体蓄热与导热，$\rho_s,c_s,k_s,T_s,s_s$ 分别为固体密度、比热、导热系数、温度和体积热源。整板当前规划采用铜，具体材料常数或温度函数待团队确认。在无接触热阻界面上，以同一个由固体指向流体的法向 $\bm n$ 表示连续条件：
\begin{equation}
T_s=T_f,\qquad -k_s\nabla T_s\cdot\bm n=-k_f\nabla T_f\cdot\bm n.\label{eq:interface}
\end{equation}
式\eqref{eq:interface}同时约束温度和热流；若后续引入接触热阻，则需改用温度跳跃条件。图\ref{fig:07}用于明确边界位置，不能仅以方程替代实际边界设置。
\begin{table}[htbp]
\centering\small
\caption{CFD工况与边界条件维护表}\label{tab:03}
\begin{tabularx}{\linewidth}{lYY}
\toprule
\rowcolor{ydzjlight} 设置 & 已读取案例或当前规划 & 待确认范围 \\
\midrule
模型 & 三维、层流、能量方程 & 代表性物性对照已用；新参数范围需复核适用性。 \\
入口 & $2\,\mathrm{m/s}$，$300\,\mathrm K$ & 当前整板规划沿用历史参数；局部最终对照另核。 \\
出口 & 表压0 Pa，回流$302\,\mathrm K$ & 当前整板规划，不能外推至全部历史案例。 \\
固体 & 铜 & \TBD{导热系数、比热、密度及其温度依赖} \\
水物性 & REFPROP拟合，280--360 K & 拟合评估压力101325 Pa；记录启用函数版本。 \\
热边界 & 空间三档映射；动态时序另设 & \TBD{连接区域热边界、非加热面条件} \\
求解类型 & 物性对照：瞬态；整板规划：稳态 & \TBD{最终局部连接对照类型与统计时间} \\
数值格式 & 物性对照：SIMPLE，一阶隐式时间格式 & \TBD{空间格式、松弛因子及统一生产配置} \\
残差目标 & 物性对照：连续性/速度$10^{-3}$，能量$10^{-6}$ & 每步最多20次内迭代；生产样本准入须另确认。 \\
重力/操作压力 & 代表性案例关闭重力；101325 Pa & \TBD{逐案例核对，禁止以残差单独判稳态} \\
\bottomrule
\end{tabularx}
\end{table}

\placeholderfigure{fig:07}{固液计算域与边界条件}{固液分域、速度入口、压力出口、加热底面、其他壁面及固液界面；注明入口温度、回流温度、热源覆盖区域。}{33mm}

\subsubsection{网格划分及网格无关性验证}
团队已开展多轮网格检验，重点关注翼型尾部、狭窄间隙和连接区域，采用局部加密并检查网格质量。尾部剪切与回流区域对压力损失和速度梯度较敏感，单纯增加全域单元数不一定能获得相同的局部解析能力。网格质量通过是计算准备条件，网格无关性则需要比较一致物性、边界、统计时间和收敛条件下的响应量，两者不能互相替代。

本报告提取了 NACA2416\_v3 在 $0.2\,\mathrm{s}$ 的七档已有结果，见表\ref{tab:04}。其中“平均温度”沿用原数据的全域单元算术平均，并非体积加权、受热面平均或流体混合温度；“最高温度”为全域最大值；压降为进出口面积平均静压差。该数据只能按原定义解释，不作为连接结构的优势数据。图\ref{fig:08}和图\ref{fig:09}用于补入局部网格与响应曲线。
\begin{table}[htbp]
\centering\small
\caption{NACA2416\_v3七档网格在0.2 s的已有结果}\label{tab:04}
\begin{tabularx}{\linewidth}{lYYY}
\toprule
\rowcolor{ydzjlight} 单元数 & 平均温度/K & 最高温度/K & 静压差/Pa \\
\midrule
260169 & 304.668393 & 344.675487 & 2505.203605 \\
575129 & 304.411275 & 345.265735 & 2389.274100 \\
1073583 & 303.944909 & 345.153612 & 2440.765641 \\
1524876 & 305.093298 & 346.048803 & 2389.938449 \\
2943269 & 304.373551 & 345.964030 & 2355.838476 \\
3235081 & 304.300924 & 345.983833 & 2348.320348 \\
4610769 & 305.985989 & 346.618797 & 2448.714104 \\
\bottomrule
\end{tabularx}
\par\smallskip\begin{minipage}{\linewidth}\footnotesize 温度为全域单元算术平均及全域最高值，静压为进出口面积平均；非受热面指标。最细档有不同时间步历史，未通过全序列1\%联合判据。来源：evidence/mesh\_results\_02s.csv。\TBD{本报告新工况的网格验证与生产网格选择}。\end{minipage}
\end{table}

\placeholderpair{fig:08}{不同网格规模的局部网格展示}{较粗网格}{相同剖面、翼型尾缘及最小间隙；注明单元数与局部尺寸。}{加密网格}{保持相同视角与放大倍率，标出加密区域与边界层。}{28mm}
\placeholderfigure{fig:09}{网格敏感性验证曲线}{使用evidence中的原始CSV；横轴单元数，温度和压降分别绘图；注明0.2 s、统计口径与参考网格。}{32mm}

记录显示，2943269 与 3235081 单元之间的静压差相对变化较小，但 3235081 与 4610769 单元之间的静压差相对后者仍为约 $4.10\%$，不能宣称七档数据全部通过 $1\%$ 联合判据。最细档还保留了部分不同时间步历史，应作为解释限制记录。已有工作建立了网格敏感性分析和指标核验基础；后续样本准入仍需针对具体几何检查关键输出，而不是沿用一个通用单元数。绝对温度相对误差还可能掩盖温升误差，后续建议同步比较 $T-T_{\mathrm{in}}$。

\subsubsection{温度相关物性UDF与非均匀热源设置}
水物性实现已进行修正和对照，当前可用拟合 UDF 根据温度返回密度、动力黏度、导热系数和比热，比热积分与显热使用同一多项式。已读取版本声明物性拟合压力为 $101325\,\mathrm{Pa}$、材料温度范围为 $280$--$360\,\mathrm{K}$。该压力是物性评估条件，不表示流动压力场被固定。已排查修正的旧物性问题作为方法改进背景保留，不作为当前主要研究障碍。表\ref{tab:05}记录当前物性与热源配置，图\ref{fig:10}预留空间和时间热源说明。
\begin{table}[htbp]
\centering\small
\caption{当前物性模型、热源UDF与实现核对}\label{tab:05}
\begin{tabularx}{\linewidth}{lYY}
\toprule
\rowcolor{ydzjlight} 类别 & 已知配置或核对值 & 适用边界与待补项 \\
\midrule
密度 & $996.556916\,\mathrm{kg/m^3}$（300 K） & water\_rp\_density；仅当前拟合版本。 \\
比热 & $4180.636970\,\mathrm{J/(kg\cdot K)}$（300 K） & water\_rp\_cp；焓由相同比热积分。 \\
动力黏度 & $8.537430\times10^{-4}\,\mathrm{Pa\cdot s}$（300 K） & water\_rp\_viscosity。 \\
导热系数 & $0.609500\,\mathrm{W/(m\cdot K)}$（300 K） & water\_rp\_conductivity。 \\
拟合范围 & 280--360 K，101325 Pa & 物性拟合条件；超范围需扩展或更换模型。 \\
实现对照 & 共同0.137 s：加热面最高温差约$3.197\times10^{-6}$ K & 同初场、步长；不是绝对物理精度验证。 \\
静态空间热源 & 75/350/600 W/cm$^2$；70/27/3格 & 整格加热2640 W；主体面积加热2021.25 W。 \\
动态热源 & 已探索随机波动、趋势与峰值 & \TBD{最终区域热源版本、种子、时间步更新规则} \\
\bottomrule
\end{tabularx}
\end{table}

代表性对照继承同一 $0.1\,\mathrm{s}$ 初场，以 $0.001\,\mathrm{s}$ 步长推进，在共同 $0.137\,\mathrm{s}$ 比较拟合与实时 REFPROP 实现：加热面最高温度差约 $3.197\times10^{-6}\,\mathrm{K}$，面积平均压降差绝对值约 $0.009682\,\mathrm{Pa}$。这些结果支持该案例下两种物性实现的一致性，并不证明所有边界、全部物理模型或实验精度。共同前37步的接口计时比约11.50仅属于物性实现对照，不能写成代理模型加速比。

非均匀热源分为空间映射与时间演化两类。当前整板静态布局为70个低负荷、27个中等负荷和3个高负荷格，对应75、350和600 $\mathrm{W/cm^2}$。按完整 $4\,\mathrm{mm}\times4\,\mathrm{mm}$ 格加热，总输入为2640 W；若只加热 $3.5\,\mathrm{mm}\times3.5\,\mathrm{mm}$ 主体，总输入为2021.25 W。二者是不同覆盖定义，不能把后一面积与前一功率同时使用。映射代码已有中心位置、数量和总热量检查记录，当前整板几何说明进一步核对了模块与热源布局。
\placeholderpair{fig:10}{非均匀热源空间分布与时间演化}{空间映射}{100格编号、70/27/3分类及热点聚集；标注热源面积和功率。}{时间变化}{补入实际动态热源轨迹、采样周期和随机种子；与静态工况分开。}{30mm}

\subsubsection{数值仿真自动化流程}
自动化探索覆盖参数配置、几何生成、流体域处理、网格划分、求解配置、监测输出与结果归档。批量计算流程已初步建立，但不同结构的实体有效性、局部网格和求解收敛仍需检查，不能把脚本正常退出等同于样本可靠。每个算例应绑定唯一几何标识、物性版本、边界配置、网格文件及原始监测表，以支持后续复现。

图\ref{fig:11}概括自动化流程与样本准入关口。后续拟在数据集建设中增加样本准入规则：不满足残差、质量守恒、能量平衡或温度范围要求的结果应保留异常标签，重新计算或单独记录，不直接进入训练。已经完成的自动化基础为流程初步贯通，覆盖全部几何和全部工况的稳健批处理仍属于持续完善内容。
\reportfigure{fig:11}{数值仿真自动化与样本质量检查流程}{figures/diagrams/fig_11.pdf}{0.98\linewidth}

\subsubsection{流动与传热性能评价指标}
基本流动尺度采用水力直径和雷诺数：
\begin{equation}
D_h=\frac{4A_c}{P_w},\qquad Re=\frac{\rho U_bD_h}{\mu}.\label{eq:re}
\end{equation}
$A_c$ 为选定截面的流通面积，$P_w$ 为湿周，$U_b$ 为体积平均速度。复杂扰流结构的代表截面与物性参考温度必须统一，避免比较对象使用不同定义。平均换热系数和努塞尔数可定义为
\begin{equation}
h_{\mathrm{eff}}=\frac{Q_{sf}}{A_{sf}(\overline T_{sf}-T_{b,\mathrm{ref}})},\qquad Nu=\frac{h_{\mathrm{eff}}D_h}{k_{f,\mathrm{ref}}}.\label{eq:nu}
\end{equation}
$Q_{sf}$ 为固液界面实际传热量，$A_{sf}$ 为对应面积，$\overline T_{sf}$ 为界面面积平均温度，$T_{b,\mathrm{ref}}$ 为声明的流体参考温度。该式为整体等效定义，非均匀或瞬态条件下必须明确平均方式；流体出口温度宜采用质量流量加权，回流时需声明统计规则。
\begin{equation}\begin{aligned}
R_{\theta,\max}&=\frac{T_{w,\max}-T_{\mathrm{in}}}{Q_{\mathrm{in}}},&
\Delta p_s&=\overline p_{\mathrm{in},A}-\overline p_{\mathrm{out},A},\\
P_h&=\Delta p\,\dot V,& P_e&=\frac{P_h}{\eta}.
\end{aligned}\label{eq:resistance}
\end{equation}
$T_{w,\max}$ 为指定受热壁面最高温度，$Q_{\mathrm{in}}$ 为输入热功率，$\dot V$ 为体积流量，$\eta$ 为泵与驱动的综合效率。稳态 $R_{\theta,\max}$ 可用于温升比较；瞬态同一比值含有历史和蓄热影响，不能直接解释为静态热阻。面积平均静压差用于流阻对照，但存在截面速度差异或回流时，泵送评价应补充总压或机械能通量，并计入系统管路与阀门损失。
\begin{equation}
\Delta T_A=T_{w,\max}-T_{w,\min},\qquad
\sigma_{T,A}=\left[\frac{1}{A}\int_A(T_w-\overline T_{w,A})^2\,\mathrm dA\right]^{1/2}.\label{eq:uniformity}
\end{equation}
$A$ 为统一的受热面评价区域，二者分别反映极差和整体温度分散程度。表\ref{tab:06}给出指标维护模板。项目后续将优先保证指标含义与统计位置一致，再开展性能排序。
\begin{table}[htbp]
\centering\small
\caption{热工与流动性能指标定义及维护模板}\label{tab:06}
\begin{tabularx}{\linewidth}{lYY}
\toprule
\rowcolor{ydzjlight} 指标 & 统计定义 & 项目待补记录 \\
\midrule
$T_{w,\max}$ & 指定受热面最高温度/K & \TBD{监测面、时刻、实际数值} \\
$\overline T_{w,A}$ & 受热面面积平均/K & \TBD{面范围和面积加权数据} \\
$\Delta T_A,\sigma_{T,A}$ & 温度极差与面积标准差/K & \TBD{统一评价区域与数值} \\
$R_{\theta,\max}$ & 最高温升/输入功率，K/W & \TBD{功率覆盖与稳态条件} \\
$h_{\mathrm{eff}},Nu$ & 界面传热与声明流体参考温度 & \TBD{面积、参考温度与物性} \\
$\Delta p_s$ & 进出口面积平均静压差/Pa & \TBD{统计面、静压/总压区别} \\
$P_h,P_e$ & 水力功与泵电功率/W & \TBD{系统压降、流量、效率} \\
守恒误差 & 质量与能量平衡；瞬态含储存项 & \TBD{计算方式、阈值与准入记录} \\
\bottomrule
\end{tabularx}
\end{table}

\subsection{区域化冷板连接结构设计与阶段性成果（方向二）}
\label{ch:connection}
\subsubsection{区域化冷板设计思路}
方向二针对不同热负荷区域散热需求不同、模块连接可能引入额外流动损失的问题，将整体冷板拆分为可配置的基本单元，并研究连接段如何组织上游流动进入下游结构。区域化设计并不意味着单元间热量和流量完全独立：固体横向导热、上游流体温升和行级压力耦合仍会影响局部工作点。因此，局部结构适配和连接整流应作为一个组合问题评价。

项目已完成整体分区与单元设计思路，形成不同负荷类别的结构配置路径；已确定连接方案，并根据团队阶段说明取得仿真对照支持的性能优势。现阶段研究已由概念方案进入性能验证与结果整理，接近阶段性完成。本章重点展示这一实质进展，同时保留尚未归档的对照数据和适用边界。

\subsubsection{基本单元与连接结构}
基本单元按4 mm节距组织，当前整板版本将3.5 mm主体布置在格中心，行内连接长0.5 mm，两端各保留0.25 mm；条间填充与边缘处理用于满足整体外形尺寸。连接段的功能不仅是几何拼接，还包括将不同主体的出口速度分布调整为可接受的下游入口条件，降低突变接口造成的流动偏置。图\ref{fig:13}用于补入最终方案及接口尺寸。

已有连接方案库包含梭形导叶、蜂窝整流格、双排导叶及螺旋槽等版本，方案存在本身不代表均已完成性能计算。当前四条40 mm验证装配记录明确采用低阻梭形直导叶、保留六片导叶和横向轮廓，并对流向长度进行缩短。表\ref{tab:08}仅记录这一可追溯的装配版本；团队已验证连接方案与此版本的对应关系、最终成果版本号仍需确认，不能因多个目录出现“最终版”便推定为同一性能对照对象。
\begin{table}[htbp]
\centering\small
\caption{可追溯连接装配版本及最终对照参数模板}\label{tab:08}
\begin{tabularx}{\linewidth}{lYY}
\toprule
\rowcolor{ydzjlight} 参数 & 当前40 mm条带记录 & 最终性能方案待核对 \\
\midrule
连接型式 & 低阻梭形直导叶 & \TBD{最终对照版本及文件哈希} \\
导叶数量 & 6片 & \TBD{确认最终方案数量及轮廓} \\
内部/端部长 & 0.5/0.25 mm & \TBD{局部对照采用长度} \\
外壳宽/高 & 3.5/3.0 mm & \TBD{是否与最终对照一致} \\
通腔宽/高 & 3.2/2.5 mm & \TBD{导叶占据后的有效面积} \\
侧壁/上下壁 & 0.15/0.25 mm & \TBD{主体接口尺寸与过渡} \\
基准连接 & \TBD{无整流或其他基准几何} & \TBD{相同长度、截面和主体条件} \\
\bottomrule
\end{tabularx}
\end{table}

\placeholderfigure{fig:13}{已确定连接结构及单元接口}{补入最终用于性能对照的连接件；标注六片导叶或实际结构、连接长度、截面、接口；确认与当前CAD版本对应。}{33mm}

\subsubsection{三类代表性热负荷工况}
研究按低热负荷、高稳定热负荷和动态波动热负荷组织三类工作模式。低负荷区域重点考察在满足温度要求时控制流动代价；高稳定负荷区域重点考察持续热点和局部传热能力；动态负荷区域重点考察波动下的瞬态温度与后续控制接口。图\ref{fig:12}概括各模式的评价与结构需求，表\ref{tab:07}维护相应边界说明。
\begin{table}[htbp]
\centering\small
\caption{三类代表性热负荷工作模式}\label{tab:07}
\begin{tabularx}{\linewidth}{lYY}
\toprule
\rowcolor{ydzjlight} 工作模式 & 研究侧重点 & 边界与证据维护 \\
\midrule
低热负荷 & 满足温度要求并控制流动代价 & 静态布局有75 W/cm$^2$档；\TBD{实际结构与对照结果}。 \\
高稳定负荷 & 持续热点、局部热阻与连接作用 & 静态布局有600 W/cm$^2$档；\TBD{持续时间与最终工况}。 \\
动态波动负荷 & 响应、恢复和泵阀调控接口 & \TBD{波动基值、幅度、周期、随机轨迹}；不以350 W/cm$^2$代替时间模型。 \\
\bottomrule
\end{tabularx}
\end{table}

\reportfigure{fig:12}{三类代表性负荷与区域结构配置}{figures/diagrams/fig_12.pdf}{0.98\linewidth}
当前整板静态75、350、600 $\mathrm{W/cm^2}$ 三档是空间布局参数，三类工作模式则包含时间性质，二者不是完全相同的分类。尤其350 $\mathrm{W/cm^2}$ 只能代表一个负荷等级，不能证明已经覆盖动态变化。后续对照应注明热源是否恒定、变化幅度、持续时间及总功率，建立可复用工况编号。

\subsubsection{连接结构的数值验证}
依据团队提供的真实进展，连接结构已通过数值仿真对照验证性能优势。该阶段成果意味着已开展基准与改进结构的比较，而不是仅有建模设想。本次文件检查取得了连接几何与装配记录，但未定位到能够唯一对应最终方案的完整成对性能数据，因此本报告保留定性阶段结论，具体温度、压降、提升幅度及统计方法待归档补入。

正式结果整理应把比较对象控制在相同主体、相同热负荷和一致入口条件下。若比较固定流量，应报告改善带来的压降变化；若比较固定泵送功率，应允许流量不同并说明运行点求解方法。基准连接必须定义其几何和功能，避免用不同主体或不同负荷的数据归因于连接段。网格设置还应考虑连接狭窄通路，防止几何差异与局部离散差异混杂。

建议最终证据链包含：方案版本与几何图、求解配置、基准与连接方案的监测量、连接处速度与流线、温度分布和压降统计。图\ref{fig:14}为流动证据预留位置。已有局部验证与后续整板验证应使用不同案例编号。
\placeholderpair{fig:14}{连接处速度场与流线对照}{基准连接}{相同剖面、色标、入口流量；标出偏流、回流与射流位置。}{确定方案}{同一显示方式下展示流动重构；用局部截面统计支持图像解释。}{33mm}

\subsubsection{性能对照与优势分析}
表\ref{tab:09}建立基准与确定方案的最终对照模板。当前可以确认的是方案已确定、已有局部仿真支持优势；尚不能确认优势体现为哪些指标同时改善、其幅度多大、是否覆盖全部负荷与流量。填写表格时，应同时列出受热面最高温度、面积平均温度或均匀性、压降及泵送代价，对不利变化也应如实记录。
\begin{table}[htbp]
\centering\small
\caption{连接结构与基准结构的成对性能比较模板}\label{tab:09}
\begin{tabularx}{\linewidth}{lYYY}
\toprule
\rowcolor{ydzjlight} 指标 & 基准结构 & 确定方案 & 比较条件 \\
\midrule
$T_{w,\max}$/K & \TBD{数值} & \TBD{数值} & 同热源、同流量或同泵功。 \\
$\overline T_{w,A}$/K & \TBD{数值} & \TBD{数值} & 相同受热面及统计时刻。 \\
$\sigma_{T,A}$/K & \TBD{数值} & \TBD{数值} & 相同区域及面积加权。 \\
$\Delta p$/Pa & \TBD{数值} & \TBD{数值} & 声明静压、总压或能量口径。 \\
$R_\theta$/(K/W) & \TBD{数值} & \TBD{数值} & 声明稳态与总热功率。 \\
$P_h$/W & \TBD{数值} & \TBD{数值} & 流量与系统压力范围一致。 \\
优势与代价 & \TBD{原始数据编号} & \TBD{方向及幅度} & 仅根据归档数据填入，不编造百分比。 \\
\bottomrule
\end{tabularx}
\par\smallskip\begin{minipage}{\linewidth}\footnotesize 已有局部仿真支持优势为团队提供的阶段事实；本表具体量化尚待原始对照归档。\end{minipage}
\end{table}

\placeholderfigure{fig:15}{连接结构压降对照}{补入相同流量或相同泵功条件下基准与连接方案数据；注明静压/总压、统计截面、误差及工况编号。}{28mm}
\placeholderpair{fig:16}{连接结构温度场与换热效果对照}{温度场}{基准与方案同一受热面、相同色标和时刻；标出最高温位置。}{定量换热指标}{补入温升、热阻或Nu；声明面积与流体参考温度，避免仅凭云图判断。}{31mm}
图\ref{fig:15}和图\ref{fig:16}对应压力代价与热响应两个维度。优势分析应回答：改善是否出现在局部连接附近，是否延续到下游单元，是否随入口流量变化，以及是否以额外阻力为代价。如果不同指标呈现权衡，可以表述为特定工况下的综合适用性，而不能用一个未经声明权重的综合分数掩盖取舍。
\pending{需补入确定方案的成对原始结果、工况编号、温度与压力监测位置、比较约束及误差；不得填写未经提供的提升百分比。}

\subsubsection{流动传热机理讨论}
连接段可能通过改变截面收缩与扩张过程、引导横向速度分量和调整下游入口分布影响传热。若导叶使入口偏流减弱，下游部分通道可能获得更均衡供液；若强化局部混合，则可能改变温度梯度和边界层发展。这些是需要由速度、流量分配与温度证据验证的机理假设，不能因为设置了导叶便认定所有机制已经成立。

建议沿连接前、连接中和连接后三个位置提取截面速度分布与局部压力，比较不同支路流量及下游受热面温度。若观察到回流减弱，需要同时确认涡区大小、持续位置与流量通行情况，而不只展示某一条流线。若温度改善发生于下游较远处，应讨论流体温升、固体导热和入口重新发展共同作用，避免将全部变化归因于单一局部现象。

由于泵送功率取决于流量和系统阻力，局部压降降低不必然等于系统电能降低；不同进出口速度分布还可能改变静压与总压差的对应关系。后续机理解释将结合已建立的多口径压力核验方法，选择与研究问题匹配的压力评价。图\ref{fig:17}概括几何、流动、热响应与能耗之间的分析链条，图中的因果箭头应由对照结果支持。
\reportfigure{fig:17}{连接结构性能优势的综合分析}{figures/diagrams/fig_17.pdf}{0.98\linewidth}

\subsubsection{当前完成情况与成果总结}
方向二已完成基本单元划分、负荷分类、连接方案确定及局部数值对照，具有阶段性实质成果，相关研究接近阶段性完成。当前工作主要是确认最终版本、整理可追溯对照、分析关键参数并完成机制说明，而不是重新将研究描述为只有结构构想。现有40 mm条带和40 mm整板装配为扩展验证提供几何基础，整板性能计算与系统级结论仍在后续范围内。

阶段性论文和相关成果准备将以已经归档的局部验证为依据，先完成数据和图表，再形成结果讨论。论文撰写准备不等于投稿、录用或发表；本报告不设置任何未经确认的发表成果。下一步优先形成能由指导教师独立核查的几何、边界、原始数据与图表组合，使阶段优势从团队进展说明转化为完整报告证据。
\subsection{非均匀动态热负荷下的协同调控研究（方向三）}
\label{ch:control}
\subsubsection{动态热负荷问题定义}
方向三关注热负荷在空间和时间上变化时，如何通过有限数量的泵阀执行器调节冷却供给，使热点温度受到约束，同时改善温度分布并控制泵送代价。与稳态结构比较不同，动态控制的评价对象是一段完整响应过程，包括扰动开始、控制动作、温度过渡及恢复，而不是单一时刻的最低温度。

可将受热面划分为区域 $A_i$，用 $q_i''(t)$ 描述第 $i$ 区热流密度。项目采用十行单向流路的初步规划，行级控制对应的是一组区域的集体供液，不能独立指定100个单元的流量。沿流向的流体升温、横向固体导热以及共用供回水压力使控制对象具有耦合性。相邻区域的温度响应可能受到同一控制动作影响，这正是协同控制需要处理的问题。

当前已形成基本架构并完成最小闭环初步验证，尚未形成完整的多场景控制性能数据。后续研究需首先界定允许最高温度、反馈区域、采样周期和泵阀限制，再比较控制算法。未知温度约束和实际硬件特性均保留待确认，不预设能够满足的性能指标。

\subsubsection{双泵与行级阀门控制架构}
系统基本方案为双泵与行级出口阀门协同调控，利用区域温度反馈决定泵工作状态和各行阀门开度。泵主要影响供液压力或总流量，阀门主要改变各行水力阻力；两者通过公共管路作用于全部行，不能简单解释为泵负责温度、阀门负责均匀性而完全独立。图\ref{fig:18}示意系统架构与温度反馈链，表\ref{tab:10}给出输入、状态与输出维护模板。
\reportfigure{fig:18}{双泵与行级出口阀门协同控制架构}{figures/diagrams/fig_18.pdf}{0.98\linewidth}
\begin{table}[htbp]
\centering\small
\caption{动态系统的输入、状态、测量及控制输出}\label{tab:10}
\begin{tabularx}{\linewidth}{lYY}
\toprule
\rowcolor{ydzjlight} 类别 & 候选量与物理意义 & 当前状态或待确认项 \\
\midrule
扰动输入 & $q_i''(t)$、入口温度 & 动态热源已探索；\TBD{实际区域时序}。 \\
热状态 & $T_i,T_{f,i}$，区域与流体温度 & \TBD{状态维数、集总假设与传感器} \\
水力状态 & $\dot V_i,\Delta p$，行流量与供液压力 & \TBD{测量方式、动态/代数处理} \\
测量输出 & 区域最高温、代表点或平均温度 & 已有反馈思路；\TBD{闭环实际使用定义}。 \\
控制输出 & 两泵转速/流量指令、行级出口阀开度 & 基本架构已提出；\TBD{回路、阀数量和动作限制}。 \\
约束 & 温度上限、最小流量、动作速率 & \TBD{真实阈值与执行器曲线} \\
\bottomrule
\end{tabularx}
\end{table}

当前明确的是双泵和行级阀门的基本框架，两台泵的串并联关系、实际回路分组、阀门数量与可调范围、温度测量方式仍需结合系统方案确认。十行规划通常需要逐行控制接口，但未提供执行器采购或标定记录，因此报告不将“十个实物阀门已安装”作为成果。反馈温度可取区域最高值或代表点，但选择不同统计方式会改变控制动作，需统一记录。

建议后续采用分层协调思路：上层根据最热区域及总需求调节泵供给，下层依据行间需求差异分配阻力，同时施加流量下限和动作速率约束。这是可研究的实现方式，尚不作为已完成算法。需要特别防止阀门同时节流导致泵供给与行流量目标冲突，或某行开启使其他行流量下降。

\subsubsection{动态非均匀热源建模}
一般区域热源可表示为
\begin{equation}
q''(\bm r,t)=\sum_{i=1}^{100}\chi_i(\bm r)\left[q_{i,0}''+a_i s_i(t)\right],\qquad
Q_{\mathrm{in}}(t)=\sum_i A_i q_i''(t).\label{eq:heatmap}
\end{equation}
$\bm r$ 为受热面坐标，$\chi_i$ 为区域指示函数，$q_{i,0}''$ 为基值，$a_i$ 为变化幅度，$s_i(t)$ 为无量纲时间信号。各信号需满足热流非负与目标范围要求。阶跃、脉冲、周期波动和随机轨迹用于分别观察响应速度、恢复能力、周期跟踪与扰动适应性；式\eqref{eq:heatmap}是后续统一建模表达，不代表100区随机独立热源已全部实现。

现有热源 UDF 已探索升温趋势、随机波动与短时峰值。其中读取到的一个历史文件在同一次调用中向同一壁面所有网格面赋相同热流值，因而主要描述时间变化，不能仅凭该文件认定已经实现空间独立波动。团队提供的动态非均匀研究进展作为阶段状态记录，最终需补入实际使用的区域热源文件、版本与输出轨迹，以证明空间差异和时间变化同时存在。

随机热源用于算法比较时，还需要可复现性。同一随机种子本身不足以保证跨计算模式完全一致，因为 UDF 调用次数和并行执行可能不同。后续宜采用预生成热源时间表或按物理时刻索引的确定性映射，将扰动更新绑定到时间步而非不受控的函数调用。控制算法之间应复用同一轨迹和初态，避免以不同扰动比较优劣。这属于热源模型深化任务，并非否定已完成的可行性探索。图\ref{fig:19}将以同一时间轴对应真实负荷与温度。
\placeholderpair{fig:19}{动态热源轨迹与区域温度响应}{热源轨迹}{补入至少两区域或两行的真实热源时间序列，注明种子、单位和更新时间。}{温度响应}{与左图同一时间轴；标记扰动、控制动作及热点位置变化。}{31mm}

\subsubsection{初步数学模型与最小闭环验证}
项目已完成最小闭环的初步可行性验证，其意义是确认温度测量、控制计算、流量相关动作和温度反馈能够形成基本闭环。由于本次未获得闭环日志、程序版本和对照轨迹，报告不填写超调、响应时间或节能幅度。图\ref{fig:20}与表\ref{tab:11}用于完整补入这一阶段证据，记录验证是在数值耦合平台还是实物平台开展、启用多少执行器以及反馈如何更新。
\placeholderpair{fig:20}{最小闭环控制的初步验证结果}{开环与闭环温度}{补入同一热源与初态下的实际曲线；声明这是初步可行性还是性能比较。}{控制与流量记录}{补入泵指令、阀门开度、行流量及采样时刻；不预画示意结果曲线。}{32mm}
\begin{table}[htbp]
\centering\small
\caption{最小闭环验证记录与性能评价模板}\label{tab:11}
\begin{tabularx}{\linewidth}{lYY}
\toprule
\rowcolor{ydzjlight} 条目 & 已有状态 & 需补原始指标或证据 \\
\midrule
平台与链条 & 已完成初步可行性验证（团队说明） & \TBD{数值/实物平台、程序版本与日志} \\
扰动与初态 & 已开展动态热负荷模拟 & \TBD{实际热源轨迹及初始温度/流量} \\
反馈与执行器 & 双泵/行级阀门路线已提出 & \TBD{最小闭环中实际启用的执行器} \\
温度响应 & 不预填效果 & \TBD{最高温、超调、恢复时间} \\
均匀性 & 不预填效果 & \TBD{受热面极差或面积标准差} \\
控制动作 & 不预填效果 & \TBD{泵指令、阀开度、饱和次数与流量} \\
能量代价 & 没有已提供的实际节能结果 & \TBD{水力功积分或实测电耗及效率} \\
对照条件 & 拟设置固定流量或开环基准 & \TBD{同热源、同初态、同约束的基准曲线} \\
\bottomrule
\end{tabularx}
\end{table}

最小闭环可以在少量状态和执行器条件下验证逻辑，但其成功不能直接推广到十行耦合回路。控制指令能否转化为实际流量、反馈周期是否足以捕捉负荷变化，以及执行器饱和后的温度是否仍受约束，均需要进一步研究。当前阶段应将“反馈链条可行”与“多场景性能优于基准”分开陈述。

基础 PID 可写为
\begin{equation}\begin{aligned}
u_i(t)&=\operatorname{sat}\left[u_{i,0}+K_{p,i}e_i(t)+K_{i,i}\int_0^t e_i(\tau)\,\mathrm d\tau+K_{d,i}\frac{\mathrm de_i}{\mathrm dt}\right],\\
e_i&=T_i-T_{i,\mathrm{set}}.
\end{aligned}\label{eq:pid}
\end{equation}
$u_i$ 为冷却动作的抽象量，$u_{i,0}$ 为基准指令，$T_i$ 与 $T_{i,\mathrm{set}}$ 为反馈与设定温度，三个增益分别对应比例、积分和微分作用，$\operatorname{sat}$ 为执行器限幅。若增大 $u_i$ 表示增强冷却，上述误差符号与正比例增益相匹配；阀开度与流量关系仍需标定。积分抗饱和、微分滤波和离散实现是后续配置，不表示当前最小闭环一定采用此完整形式。

\subsubsection{向物理机理融合模型的发展}
后续研究将由简化数学模型逐步过渡至能反映蓄热、流量和压力耦合的物理灰箱模型。对于区域集总温度，在内部温差相对可控、参数经过辨识的条件下，可建立
\begin{equation}
C_i\frac{\mathrm dT_i}{\mathrm dt}=Q_i(t)-G_i(\dot m_i)(T_i-T_{f,i})
-\sum_{j\in\mathcal N_i}G_{ij}(T_i-T_j).\label{eq:rc}
\end{equation}
$C_i$ 为区域等效热容，$G_i=1/R_i$ 为流量相关等效换热热导，$T_{f,i}$ 为流体参考温度，$G_{ij}$ 表示相邻区域导热。式\eqref{eq:rc}描述蓄热、散热和横向耦合，参数不能仅以材料常数直接替代，需要通过 CFD 或响应数据辨识。流体沿程温升明显时，应增加流体状态或分段输运描述，而不能始终使用相同入口温度。

流量分配应满足节点守恒与支路压降一致性。以共用供回液压力的并联行为例，可写
\begin{equation}\begin{aligned}
\sum_i\dot V_i&=\dot V_{\mathrm{tot}},\\
\Delta p_{\mathrm{man}}&=\Delta p_{\mathrm{ch},i}(\dot V_i,T_i)
+\Delta p_{\mathrm{v},i}(\dot V_i,\alpha_i).
\end{aligned}\label{eq:hydraulic}
\end{equation}
$\alpha_i$ 为第 $i$ 行阀门开度，$\Delta p_{\mathrm{man}}$ 为供回液压差；通道和阀门关系由结构、温度及特性曲线决定。泵特性需写为 $\Delta p_{\mathrm p}=F_{\mathrm p}(n,\dot V)$，其中 $n$ 为泵速。式\eqref{eq:hydraulic}适用于所声明的并联拓扑；两泵回路尚未确认，实际网络若不同应重新列式，不先写定串并联结构。

将区域温度、必要流体状态与压力变量组合后，可以形成非线性状态模型，并在特定工作点线性化：
\begin{equation}\begin{aligned}
\dot{\bm x}&=\bm f(\bm x,\bm u,\bm d;\bm\theta),\quad\bm y=\bm g(\bm x,\bm u),\\
\delta\dot{\bm x}&=\bm A\delta\bm x+\bm B\delta\bm u+\bm E\delta\bm d.
\end{aligned}\label{eq:state}
\end{equation}
$\bm x,\bm u,\bm d,\bm y,\bm\theta$ 分别为状态、控制、热负荷扰动、测量输出和待辨识参数；$\bm A,\bm B,\bm E$ 为工作点附近的线性矩阵。该形式是模型发展路线，当前尚未提供经辨识验证的参数矩阵。图\ref{fig:21}用于区分简化反馈、灰箱模型和后续控制比较的关系。
\reportfigure{fig:21}{由简化数学模型向物理机理融合动态模型的演进}{figures/diagrams/fig_21.pdf}{0.98\linewidth}

\subsubsection{当前研究进展与下一阶段任务}
已完成工作包括基本动态研究框架、双泵与行级阀门方案、动态热源模拟、最小闭环初步可行性验证及数学模型路线。正在推进的重点是将守恒关系和设备特性加入模型，明确反馈与执行器边界，并比较 PID 及其他适合方法。本项目不将完整深度强化学习系统作为已建成果，也没有可以报告的实际节能效果。

下一阶段先归档最小闭环证据，再建立热阻热容与水力网络的可辨识模型，使用阶跃和随机轨迹独立验证动态预测。控制比较拟设置固定流量、基础 PID 与经论证选择的物理模型方法，评价最高温度、均匀性、超调、恢复时间、累计泵送能量和动作频率。只有在热源、初态、执行器约束及统计窗口相同的条件下，才讨论相对收益。若泵效率尚未知，则报告水力功或其积分，不将其等同于实测电耗。
\subsection{物理机理融合的代理模型与计算加速研究（方向一）}
\label{ch:surrogate}
\subsubsection{代理模型研究需求的提出}
方向一源于结构与工况研究扩大后的计算需求。高保真 CFD 提供流场、温度场和边界响应，但多参数、多负荷和多次验证需要反复求解，成本包含建模、网格、求解、数据提取以及质量检查。代理模型拟在已声明的参数范围内，以少量计算替代部分重复性能预测，从而支持更密集的分析和控制模型构建。

该方向的核心问题是预测方法是否准确、数据是否足够、物理信息是否带来可验证收益，以及建库和训练成本能否在后续调用中回收。几何参数化用于形成可控样本变化和独立泛化测试，并非把寻找最优几何作为主要目标。当前已确定技术路线并建立部分 CFD 与自动化基础，尚未提供训练完成后的性能结果。

\subsubsection{代表性翼型结构与几何参数化}
已开展的代表性对象包括 NACA2416 等翼型扰流结构。当前40 mm条带装配采用的局部翼型为 NACA2416\_L450um\_98989\_A00deg，记录给出0.45 mm翼型尺度、9/8/9/8/9错列排列、每块43个扰流柱和0度迎角。该配置用于代表性建模与结构审阅，不构成性能最优证明，也不能自动作为全部样本的基准。

后续可将弦长、迎角、排布间距、柱高和翼型类别等作为候选几何变量，同时考虑入口速度或雷诺数、入口温度及热边界。候选变量并不意味着全部都必须进入模型，应通过机理分析和样本可实现性筛选，防止输入维度过高而稀释样本覆盖。参数上下界、单位及生成约束需由团队确认，不能从目录名称推定连续可行域。
\placeholderfigure{fig:22}{代表性翼型扰流结构与几何参数定义}{补入NACA2416几何、弦长、厚度、迎角、间距、柱高及排列方式；标注实际采用参数和候选参数。}{32mm}
\placeholderfigure{fig:23}{翼型尾部与局部流道网格加密}{补入真实尾缘与近壁网格、BOI范围、最小尺寸和边界层；对应具体几何版本与网格文件。}{29mm}
图\ref{fig:22}与图\ref{fig:23}对应几何定义和数值解析能力。参数改变可能导致最小间隙改变、局部接触或流体域连通性变化，因此几何生成流程需先做可行性检查，再进入批量计算。其方法价值在于建立可复现的参数—几何—网格映射，而非仅生成数量较多的模型文件。

\subsubsection{CFD数据集构建方法}
数据集以算例为基本单位，每条样本应包含输入参数、物性版本、几何版本、网格统计、边界配置、输出指标、质量标签和计算成本。历史算例可用于流程验证和探索，但不同物性、不同时间点或不同压降口径的结果不能未经统一便合并为训练标签。方向一正在构建高保真样本集，尚未确认最终有效样本数。表\ref{tab:12}为候选输入和输出建立统一字典。
\begin{table}[htbp]
\centering\small
\caption{代理模型候选输入、输出与数据标识}\label{tab:12}
\begin{tabularx}{\linewidth}{lYY}
\toprule
\rowcolor{ydzjlight} 类别 & 候选内容 & 待确认与准入要求 \\
\midrule
几何输入 & 翼型、弦长、迎角、间距、柱高 & \TBD{最终变量、范围、单位与可行性限制} \\
流动输入 & 入口速度或Re、入口温度 & \TBD{参考截面和物性温度；避免重复冗余} \\
热边界输入 & 热流等级、空间类别、必要时间参数 & \TBD{稳态任务或瞬态任务分开定义} \\
物理特征 & Re、Pr、Pe、几何比值、简化基线 & \TBD{仅由可用输入计算，消融比较} \\
温度输出 & 受热面最高温、出口混合温度 & \TBD{最终输出与统计位置} \\
性能输出 & 压降、热阻或Nu & \TBD{压力口径、面积与参考温度} \\
元数据 & 几何ID、网格、物性、守恒、计时 & 已探索自动化；\TBD{有效样本数与冻结测试集}。 \\
\bottomrule
\end{tabularx}
\end{table}

\reportfigure{fig:24}{高保真CFD样本生成与数据集划分流程}{figures/diagrams/fig_24.pdf}{0.98\linewidth}
图\ref{fig:24}概括样本生成、准入和独立测试全过程。建议先建立覆盖主要参数范围的基础样本，再根据误差或信息不足区域补点。训练、验证与测试应按几何或工况组合分组，防止同一几何的近邻样本分散到不同集合导致结果偏乐观；标准化和特征选择仅在训练数据上拟合，超参数选择仅使用验证或分组交叉验证数据。

标签的物理含义需要统一。受热面最高温度、出口混合温度、静压损失、总压损失和热阻描述不同任务，不宜混为一个未定义的“综合性能”。瞬态指标必须注明时刻或统计窗口；稳态样本需记录收敛与热平衡条件。不能将网格敏感性检验中不同单元数产生的差异当作几何差异的学习目标。

\subsubsection{物理机理分析与理论推导路线}
该方向将并行推进物理理论与物理机器学习，两条路径互相支撑而不互相替代。理论路径首先明确流动与换热的控制尺度，讨论几何长度、黏性效应、热扩散和流体输运如何影响性能；数据路径检验这些解释能否在有限样本中提高预测与泛化。相关推导正在推进，不在报告中声称已获得适用于全部结构的闭式关联式。

除 $Re$ 外，可考虑普朗特数和佩克莱数：
\begin{equation}
Pr=\frac{\mu c_p}{k_f},\qquad Pe=Re\,Pr,\qquad
t_{\mathrm{adv}}=\frac{L}{U_b},\quad t_{\mathrm{diff}}=\frac{L^2}{a},\quad a=\frac{k}{\rho c}.\label{eq:scales}
\end{equation}
$Pr$ 表示动量与热扩散的相对尺度，$Pe$ 表示对流与热扩散的相对作用；$L$ 为声明的特征长度，$a$ 为对应介质热扩散率。固体与流体应分别选取物性和长度，不能用同一 $a$ 代表整个共轭域。这些量可用于解释不同负荷与流量下的响应，但能否作为有效特征需经独立模型比较。

在稳态、无其他热损失且固液能量闭合时，冷却液输运热量满足
\begin{equation}
Q_{\mathrm{in}}=\dot m\left[h_f(T_{\mathrm{out},b})-h_f(T_{\mathrm{in}})\right].\label{eq:bulkenergy}
\end{equation}
$\dot m$ 为质量流量，$T_{\mathrm{out},b}$ 为适当定义的出口混合温度；近似恒比热时可转为 $\dot m c_p\Delta T$。瞬态则必须补充固体与流体储存以及其他边界热流。该守恒关系可为出口温升预测提供物理基线或一致性检查，但不能单独确定局部最高壁温；后者还取决于局部热阻和非均匀流动。

后续理论分析拟建立适用范围明确的基准预测，再用数据学习基准未能解释的几何和局部流动影响。若简化关联式无法覆盖扰流结构，应保留其偏差而非强行把所有样本压到统一曲线。图\ref{fig:25}用于呈现两条路径之间的双向验证。
\reportfigure{fig:25}{物理推导与物理机器学习的双路径框架}{figures/diagrams/fig_25.pdf}{0.98\linewidth}

\subsubsection{物理与机器学习融合的模型框架}
基础模型拟从 XGBoost 等表格回归方法开始，建立不含额外物理信息的比较基线。XGBoost 的树提升算法体系可作为方法来源\cite{chen2016}，但其在其他任务中的表现不能直接证明本项目精度。物理信息融合可优先比较输入特征增强、简化物理基线加残差学习，以及适用于可微模型的守恒约束等方案。物理约束学习文献提供了将方程信息引入学习过程的参考\cite{raissi2017}，具体方案仍需依据本项目输出形式和数据成本选择。

对于具有可靠简化基线的输出，可采用
\begin{equation}
\widehat y(\bm z)=y_{\mathrm{phys}}(\bm z)+r_{\boldsymbol\vartheta}(\bm z,\bm\phi(\bm z)).\label{eq:residualmodel}
\end{equation}
$\bm z$ 为原始输入，$\bm\phi$ 为只依赖可用输入的物理特征，$y_{\mathrm{phys}}$ 为物理基线，$r_{\boldsymbol\vartheta}$ 为学习残差。模型输入不能包含预测时必须额外执行 CFD 才能得到的量，否则会削弱计算加速意义并造成标签信息泄漏。对最高壁温建立基线时还需说明局部热阻近似，不能直接套用仅适用于出口混合温度的能量式。

对于可微且可表达物理约束的候选模型，可研究
\begin{equation}
\mathcal L=\mathcal L_{\mathrm{data}}+\lambda_E\mathcal L_E+\lambda_B\mathcal L_B.\label{eq:loss}
\end{equation}
$\mathcal L_{\mathrm{data}}$ 为数据误差，$\mathcal L_E$ 为可计算的守恒残差，$\mathcal L_B$ 为边界或非负性等一致性惩罚，权重 $\lambda_E,\lambda_B$ 需标定并统一无量纲尺度。表格标量回归不必具备完整流场的空间导数，XGBoost 也不能直接照搬常规神经网络的 PDE 自动微分残差。物理特征增强与残差学习更适合作为树模型的初步融合路径；方程约束网络则属于需另行论证的数据与算法路线。
\begin{table}[htbp]
\centering\small
\caption{拟比较代理模型算法与物理融合路径}\label{tab:13}
\begin{tabularx}{\linewidth}{lYY}
\toprule
\rowcolor{ydzjlight} 模型路线 & 方法与比较目的 & 阶段状态与待补结果 \\
\midrule
基础回归 & XGBoost等；建立无额外物理信息基线 & 研究路线已确定；\TBD{训练与测试误差}。 \\
物理特征增强 & 加入无量纲数与可计算机理特征 & 拟比较；\TBD{消融、样本效率与泛化}。 \\
物理基线+残差 & 学习简化物理模型的剩余偏差 & 拟研究；\TBD{基线适用域与偏差}。 \\
可微约束模型 & 数据项与守恒/边界约束联合训练 & 备选路线；\TBD{输出可微性、损失与权重}。 \\
简单参考模型 & 线性/低阶回归或可验证关联式 & 建议对照；\TBD{与复杂方法公平比较}。 \\
\bottomrule
\end{tabularx}
\par\smallskip\begin{minipage}{\linewidth}\footnotesize 所有算法结果均待产生；不将技术路线写为训练完成。方程约束网络不等同于对XGBoost直接加入PDE自动微分残差。\end{minipage}
\end{table}

\reportfigure{fig:26}{物理融合代理模型架构与预测任务}{figures/diagrams/fig_26.pdf}{0.98\linewidth}
图\ref{fig:26}示意拟研究的候选模型架构，表\ref{tab:13}维护拟比较算法与结果状态。在相同训练集和测试集上比较基础模型、物理特征增强、残差模型与其他可行方法，才能判断物理信息是否有效。增加特征数量本身不属于融合有效性的证据，应使用消融分析分别移除物理特征或基线，记录误差与样本需求的变化。

\subsubsection{代理模型精度、泛化及计算效率评价}
拟以独立测试集误差、分组泛化、样本量学习曲线和计算成本共同评价。对于 $N$ 个参考输出 $y_j$ 和预测 $\widehat y_j$，可定义
\begin{equation}
\operatorname{MAE}=\frac1N\sum_{j=1}^N|\widehat y_j-y_j|,\qquad
\operatorname{RMSE}=\sqrt{\frac1N\sum_{j=1}^N(\widehat y_j-y_j)^2}.\label{eq:errors}
\end{equation}
二者与输出量具有相同单位，MAE反映平均偏差，RMSE对较大偏差更加敏感。相对误差与决定系数可作为补充：
\begin{equation}
\varepsilon_{j}=\frac{|\widehat y_j-y_j|}{\max(|y_j|,\epsilon_y)},\qquad
R^2=1-\frac{\sum_j(\widehat y_j-y_j)^2}{\sum_j(y_j-\overline y)^2}.\label{eq:relative}
\end{equation}
$\epsilon_y$ 为具有输出单位的预声明防零阈值，$\overline y$ 为测试集均值；分母为零时 $R^2$ 不适用。对于温度，建议同时报告 Kelvin 绝对误差和温升误差，不能仅用绝对温度的相对百分比展示精度。最高温预测还应检查低估尾部，以识别热约束风险。
\placeholderpair{fig:27}{代理模型预测精度与泛化比较（待获得结果）}{预测与参考}{后续填入独立测试散点及误差分布；目前无训练结果，不预填精度。}{分组与样本效率}{后续填入未见几何/流量/热负荷误差和学习曲线；与插值结果区分。}{30mm}
图\ref{fig:27}预留独立预测和泛化结果，目前尚无可填入的训练精度。泛化评价分为已覆盖范围内的插值、未见几何组合、未见流量或热边界，以及超出训练范围的外推。外推任务需单独声明，不能将随机划分的低误差等同于任意参数泛化。样本效率应在相同数据子集和独立测试集上比较多个训练规模，并记录抽样变动；精度目标和可接受误差由团队依据用途确认。

单次调用加速比与累计成本应同时报告：
\begin{equation}\begin{aligned}
S_{\mathrm{call}}&=\frac{t_{\mathrm{CFD}}}{t_{\mathrm{infer}}},\\
C_{\mathrm{sur}}(M)&=C_{\mathrm{data}}+C_{\mathrm{train}}+M t_{\mathrm{infer}},\\
C_{\mathrm{direct}}(M)&=M t_{\mathrm{CFD}}.
\end{aligned}\label{eq:cost}
\end{equation}
$M$ 为目标预测次数，$C_{\mathrm{data}}$ 与 $C_{\mathrm{train}}$ 分别为建库与训练调参的计算时间成本，$t_{\mathrm{infer}}$ 为包含必要特征处理的单次预测成本，$t_{\mathrm{CFD}}$ 为同硬件口径下的单次参考求解成本。若 $t_{\mathrm{CFD}}>t_{\mathrm{infer}}$，累计时间收益的简化回收阈值为
\begin{equation}
M>\frac{C_{\mathrm{data}}+C_{\mathrm{train}}}{t_{\mathrm{CFD}}-t_{\mathrm{infer}}}.\label{eq:breakeven}
\end{equation}
式\eqref{eq:breakeven}未计入维护与新范围补样，应在成本表中另列。并行计算应区分墙钟时间与核时，冷启动与批量预测应分别记录。已有物性接口11.50倍计时对照与这里的代理调用加速不是同一对象。表\ref{tab:14}和图\ref{fig:28}用于记录后续实测成本，当前保持占位。
\begin{table}[htbp]
\centering\small
\caption{CFD与代理模型计算开销记录模板}\label{tab:14}
\begin{tabularx}{\linewidth}{lYYY}
\toprule
\rowcolor{ydzjlight} 成本项 & 高保真CFD & 代理方法 & 计时口径 \\
\midrule
单次调用 & \TBD{求解及必要处理时间} & \TBD{特征处理与推理时间} & 同硬件，声明批量大小与冷启动。 \\
几何与网格 & \TBD{时间与核时} & \TBD{是否需要或可复用} & 与求解成本区分。 \\
样本建库 & \TBD{有效/失败样本成本} & \TBD{总建库成本} & 包括质量检查与必要重算。 \\
训练与调参 & 不适用 & \TBD{训练、验证与调参成本} & 记录数据规模与搜索次数。 \\
累计调用 & $Mt_{\mathrm{CFD}}$ & $C_{\mathrm{data}}+C_{\mathrm{train}}+Mt_{\mathrm{infer}}$ & 墙钟时间或核时采用同一口径。 \\
加速与回收 & \TBD{参考成本} & \TBD{调用加速及回收次数} & 当前没有实测代理加速结果。 \\
\bottomrule
\end{tabularx}
\end{table}

\placeholderpair{fig:28}{CFD与代理模型计算开销比较（待获得结果）}{单次调用}{后续填入同硬件CFD与推理计时，注明处理流程、重复次数与离散度。}{累计成本}{后续填入建库、训练及随调用量增长的成本和回收点；目前无实测加速结果。}{30mm}

\subsubsection{当前研究进展与后续计划}
已开展代表性翼型建模与仿真、参数化自动建模探索、初步批量 CFD 流程、多轮网格检验、温度相关物性修正，以及边界条件和评价口径研究。这些是代理模型的数据与数值基础，不能据此宣布模型已经训练完成。当前正在明确输入参数、建立合格样本集、推进理论尺度分析和 XGBoost 等基础方法研究。

后续按“统一标签—建立基础模型—加入物理信息—消融比较—泛化与成本验证”的顺序推进，优先限定可完成且可检验的标量任务。先得到可靠的基准，再评价物理信息能否降低误差或样本需求，最后判断累计成本是否支持实际重复预测。若样本范围不足或物理基线存在系统偏差，应收缩用途或补充数据，而不是提前扩大模型适用声明。几何优化可以作为远期应用讨论，但不作为本方向中期核心成果。
\section{问题与对策}
\subsection{遇到的主要问题}

\subsubsection{局部连接作用能否在不同工况及整板中保持}
连接改变入口分布，其收益可能依赖上游主体、入口流量和下游结构。需要识别改善主要来自阻力变化还是换热分布变化，明确局部结论推广到多单元串联时的条件。固体导热、沿程流体升温和接口累积损失可能削弱或改变局部优势，因此系统验证不宜简单放大局部改善比例。

\subsubsection{动态模型需要保留多少物理信息}
热阻热容模型计算轻便，但可能难以描述流体沿程温升、负荷快速波动和泵阀耦合；完整 CFD 又不利于反复控制分析。需要在可辨识性、计算代价和响应精度之间确定合理模型层级。拟先比较分区集总模型，再按误差增加流体状态或输运时延，避免一开始引入过多不可辨识参数。

\subsubsection{物理融合是否能够改善数据效率与泛化}
物理特征、物理基线与约束学习具有不同适用条件，加入物理信息也可能因简化假设不当带来偏差。需要用同一数据划分和消融比较验证其作用，研究在少样本、未见几何和负荷变化情况下是否仍有收益。仅得到训练集拟合或随机测试精度不足以回答这一问题。

\subsubsection{系统性能如何统一评价}
最高温、温度均匀性与泵送代价之间存在权衡，不同压力口径和不同热源覆盖面积会改变结论。需统一控制约束、监测位置、统计时段和成本口径。已修正的问题归入平台改进记录；当前核心科学问题集中于机制、模型层级、泛化和协同评价，而非重复罗列已排查的软件操作问题。

\subsection{已尝试的方案}
\begin{enumerate}
\item 数值可靠性方面，已有多轮网格敏感性检验、压力统计核验与水物性实现对照；这些核对明确了误差来源和统计差异，但未替代新工况的生产网格验证。
\item 动态建模方面，已开展随机热源模拟并实现最小闭环的初步验证；简化反馈可行性尚不足以回答多行水力耦合、输运时延与能耗收益问题。
\item 重复计算方面，已探索参数化与自动化流程，并提出物理融合代理路线；该路线仍需合格数据和独立测试，尚无训练完成的性能结论。
\end{enumerate}
\subsection{下一步对策}
\begin{enumerate}
\item 固定几何、物性、初场、时间步与统计口径，归档局部连接成对原始结果，并复核网格与守恒准入条件。
\item 补齐最小闭环的热源、温度、流量和执行器日志，先辨识可验证的分区热阻热容及水力参数，再按独立响应误差增加模型复杂度。
\item 统一数据划分，依次比较基础模型、物理特征与残差模型，用消融、分组泛化与累计成本检验物理融合效果。
\end{enumerate}
\section{阶段性结论}
\subsection{可以确定的}
本项目围绕时空非均匀热负荷下微通道冷板的结构适配与动态运行开展研究。初期以区域化结构及连接整流、流量分配与协同控制为两条路径；随着局部结构和工况仿真的扩展，高保真计算需求进一步催生物理机理融合的代理模型与计算加速研究。当前三个方向构成由工程问题逐步深化形成的研究体系，其发展顺序是本报告叙事的基础。

中期已有成果以方向二的连接方案确定和局部数值验证最为集中，相关工作接近阶段性完成，正在整理最终性能与机制证据。方向三已建立基本架构、开展动态热源模拟并完成最小闭环初步可行性验证，正在由简化反馈模型向能量守恒、热阻热容和水力特性融合的模型发展。方向一已有代表性翼型数值研究、网格检验、物性修正和自动化探索，正在建设数据集与研究基础算法，尚未形成独立测试支持的模型精度或加速结论。

现有资料还明确了局部模型、条带装配与整板验证之间的边界。当前整板几何和热源映射已提供后续系统研究对象，而整板流热性能、完整多行控制对比及代理模型协同使用仍需进一步验证。网格敏感性与物性实现对照的已有结果为方法规范提供依据，也提示后续必须坚持统计口径、适用范围与误差来源的明确记录。

下一阶段将优先完成连接结构证据闭合，补齐最小闭环原始记录并深化可辨识动态模型，同时以统一高保真数据评价物理融合方法的精度、样本需求与累计计算成本。在共享几何和负荷场景下，逐步比较温度响应、温度均匀性和泵送代价，形成能够追溯到原始算例与控制记录的系统结论。项目的后续发展以证据和可验证里程碑为依据，允许根据实际时间与数据范围收缩任务，保持中期报告与真实进展一致。
\subsection{尚不能确定的}
\begin{enumerate}
\item 局部连接优势的最终幅度及其整板适用性：缺少已归档的最终成对指标与统一对象上的系统求解。
\item 动态控制的温度改善与泵送代价：缺少完整控制日志、可辨识参数及相同扰动下的对照评价。
\item 代理模型的独立预测精度、泛化及累计加速收益：数据集和算法仍在推进，需独立测试与同硬件计时。
\end{enumerate}
\section{下月计划}
\label{ch:plan}
\subsection{下一阶段研究任务与里程碑}
此处承接原中期报告的下一阶段安排。月度起止日期尚未确认，故保留 M1--M4 的依赖顺序和日期待补项，不将全部里程碑承诺为一个月内完成。

下一阶段以依赖关系组织任务，具体日期在团队确认结题节点后填写，表\ref{tab:16}中的分工为建议岗位而非实际历史贡献。首先完成方向二的最终几何版本与成对结果归档，建立可复核的局部证据包；同步归档最小闭环的热源、控制和温度序列。其次统一 CFD 参数、指标与样本准入，构建基础数据集与动态参数辨识数据。
\begingroup\small
\setlength{\tabcolsep}{4pt}
\begin{longtable}{>{\raggedright\arraybackslash}p{0.09\linewidth}>{\raggedright\arraybackslash}p{0.21\linewidth}>{\raggedright\arraybackslash}p{0.28\linewidth}>{\raggedright\arraybackslash}p{0.18\linewidth}>{\raggedright\arraybackslash}p{0.15\linewidth}}
\caption{下一阶段任务、里程碑与建议分工}\label{tab:16}\\
\toprule \rowcolor{ydzjlight} 阶段 & 任务 & 可检查交付与判据 & 建议岗位（待确认） & 时间 \\\midrule
\endfirsthead
\multicolumn{5}{l}{表\thetable\ 续：下一阶段研究计划}\\
\toprule \rowcolor{ydzjlight} 阶段 & 任务 & 可检查交付与判据 & 建议岗位（待确认） & 时间 \\\midrule
\endhead
\midrule\multicolumn{5}{r}{续下页}\\\endfoot
\bottomrule\endlastfoot
M1 & 连接成果证据闭合 & 最终几何、统一工况、成对原始数据、图表与机理说明；\TBD{验收阈值} & 成员1主责；成员2复核 & \TBD{日期}\\
M1 & 最小闭环归档 & 热源、反馈、执行器、温度和流量时序；复核最小闭环实际链条 & 成员3主责；成员4复核 & \TBD{日期}\\
M2 & 数值与样本准入 & 参数字典、网格检查、物性、守恒与收敛标签；\TBD{有效样本数} & 成员2主责；成员1复核 & \TBD{日期}\\
M2 & 物理动态模型辨识 & 热阻热容及水力网络、参数辨识与独立动态响应；\TBD{误差阈值} & 成员3主责；成员2支撑 & \TBD{日期}\\
M3 & 基础与物理融合模型 & 统一划分、基础算法、融合与消融、样本效率和泛化测试 & 成员4主责；成员2支撑 & \TBD{日期}\\
M3 & 控制算法比较 & 同扰动/初态/约束；温度、均匀性、动作与能量统计 & 成员3主责；成员1支撑 & \TBD{日期}\\
M4 & 系统级协同验证 & 同一整板对象的参考CFD、动态场景、误差与成本评价；\TBD{可完成范围} & 四人共同；负责人统筹 & \TBD{日期}\\
M4 & 报告与成果整理 & 来源索引、图表定稿、阶段性论文稿及报告终审；不预设发表 & 四人共同；指导教师审核 & \TBD{日期}\\
\end{longtable}
\endgroup
\noindent\footnotesize 注：成员1--4为建议岗位标识，并非实际姓名或贡献证明；具体安排由团队确认。\normalsize

在此基础上，控制方向完成灰箱模型的独立响应验证，代理方向完成基础算法与物理融合方案的公平比较。最后开展共享对象上的系统级验证。每个里程碑应以可检查的文件、图表和判据为交付，而不是“继续优化”一类无法验收的描述。如果阶段时间不足，可优先完成连接成果收尾、最小闭环证据及一个清晰限定的代理任务，再删减超出范围的扩展内容。图\ref{fig:30}概括任务依赖和相对里程碑；实际日期待团队确认。
\reportfigure{fig:30}{下一阶段总体路线与里程碑}{figures/diagrams/fig_30.pdf}{0.98\linewidth}

\subsection{三个方向的协同与系统级验证路线}

系统验证拟统一冷板几何、热源布局与物性版本，依次复核局部连接、建立多单元或整板 CFD 参考、比较动态模型与控制方法，再评价代理模型替代重复计算的价值。每一层均需独立验证，不能直接外推上一层结论。

方向二提供结构与性能数据，方向三提供运行范围和动态需求，方向一在误差评价通过后提供快速预测。整板70/27/3静态布局可作为首个对象，动态场景再引入阶跃、脉冲和区域迁移。比较采用相同初态、热源与执行器约束；代理模型接入时须检查热点低估并保留关键场景 CFD 复核。模型联动仍属计划，最终报告场景适用性与性能权衡。
\FloatBarrier
\begin{thebibliography}{9}
\bibitem{tuckerman1981} Tuckerman D B, Pease R F W. High-performance heat sinking for VLSI[J]. IEEE Electron Device Letters, 1981, 2(5): 126--129. DOI: 10.1109/EDL.1981.25367.
\bibitem{chen2016} Chen T, Guestrin C. XGBoost: A Scalable Tree Boosting System[EB/OL]. arXiv:1603.02754, 2016. DOI: 10.48550/arXiv.1603.02754. 预印本；用于基础算法来源。
\bibitem{raissi2017} Raissi M, Perdikaris P, Karniadakis G E. Physics Informed Deep Learning (Part I): Data-driven Solutions of Nonlinear Partial Differential Equations[EB/OL]. arXiv:1711.10561, 2017. DOI: 10.48550/arXiv.1711.10561. 预印本；用于物理约束学习的方法背景。
\end{thebibliography}
\end{document}
